% Author: Adam | Module: AI Job Interview Simulator | Project: SmartCareer AI
\section{Cahier des charges}

\subsection{Exigences fonctionnelles}

\begin{description}[leftmargin=1.4em,itemsep=3pt,topsep=3pt,style=nextline]
  \item[F1 --- Génération de questions adaptées au poste.] L'utilisateur décrit
    un poste (intitulé obligatoire, entreprise et description d'offre
    facultatives), choisit un niveau parmi quatre (\const{junior},
    \const{inter}, \const{senior}, \const{lead}), un type de questions
    (\const{technical}, \const{hr}, \const{mixed}), un nombre de questions
    entre 3 et 15 et une langue. Le module renvoie exactement ce nombre de
    questions, chacune avec une catégorie, un axe (\const{focus}) et 2 à 4
    \emph{éléments de réponse attendus} qui serviront de barème à
    l'évaluation~(F2).
  \item[F2 --- Évaluation notée sur 10.] Pour une réponse donnée, le module
    renvoie une note globale décimale entre 0 et 10 et quatre critères
    indépendants notés de la même façon : \emph{pertinence par rapport à la
    question posée} (\const{relevance}), \emph{clarté} (\const{clarity}),
    \emph{structure} (\const{structure}), \emph{justesse technique}
    (\const{technical\_accuracy}). S'y ajoutent 1 à 4 forces, 1 à 4 faiblesses
    citant un passage de la réponse, et un verdict de 2 à 3 phrases.
    L'évaluation d'une réponse ne bloque pas la rédaction de la suivante : elle
    est lancée en arrière-plan et séquencée (section~\ref{sec:seq}).
  \item[F3 --- Réponse améliorée, à la demande.] Le bouton
    «~Proposer une meilleure réponse~» n'est jamais appelé automatiquement :
    chaque analyse coûte une génération. Le module renvoie une réponse réécrite,
    décomposée en STAR (Situation, Tâche, Action, Résultat) pour le
    comportemental et en démonstration de raisonnement pour le technique, plus
    2 à 3 conseils actionnables pour l'oral.
  \item[F4 --- Historique et synthèse.] Une session apparaît dans l'historique
    de l'utilisateur avec son poste, sa date, son nombre de réponses analysées
    et sa moyenne. Elle est consultable en lecture seule, et supprimable avec
    confirmation ; la suppression est bornée au propriétaire. Quand toutes les
    questions ont reçu une réponse, la session passe automatiquement à l'état
    \const{finished} et une synthèse déterministe (moyenne, meilleure et moins
    bonne question, conseils de clôture) est figée en base.
  \item[F5 --- «~Utiliser mon CV~».] L'utilisateur peut adosser les questions à
    un CV déjà généré par le Module~1, choisi dans une liste alimentée par
    l'endpoint \route{GET /history} du Module~1. Le CV est alors lu, anonymisé,
    condensé, et l'écart entre l'offre et le profil est calculé sans IA pour
    orienter les questions.
  \item[F6 --- Bilingue français / anglais.] La langue choisie s'applique aux
    questions comme à tous les textes d'analyse.
  \item[F7 --- Diagnostic.] Une route d'état indique, avant de lancer un
    entretien, si le modèle est joignable et si l'index documentaire est chargé,
    avec le nombre de fragments et le nom du modèle de plongement.
  \item[F8 --- Libération de la VRAM.] La fin de page (rechargement, fermeture
    d'onglet, navigation interne vers l'éditeur de CV) déclenche une libération
    explicite du modèle pour rendre la carte au Module~1.
\end{description}

\subsection{Exigences non fonctionnelles}

\begin{description}[leftmargin=1.4em,itemsep=3pt,topsep=3pt,style=nextline]
  \item[N1 --- Confidentialité.] Aucune donnée personnelle brute n'est envoyée
    au modèle, stockée dans une table du module ou écrite dans un journal.
    Le CV est anonymisé avant toute utilisation ; les textes libres saisis sont
    nettoyés des adresses électroniques, téléphones et URL.
    Ce point est vérifié par des tests dédiés (section~\ref{sec:tests}).
  \item[N2 --- Local et hors ligne.] Aucune API payante, aucun appel réseau
    sortant autre que \route{localhost}. Le module fonctionne câble débranché,
    une fois les poids du modèle et de l'embeddeur en cache local.
  \item[N3 --- Robustesse.] Toute panne (modèle absent, Ollama injoignable,
    timeout, quota, sortie non conforme) produit un message d'erreur français
    actionnable et un code HTTP cohérent. Aucune route ne peut renvoyer un
    500~nu avec une trace de pile.
  \item[N4 --- Cible matérielle.] Un poste de développement ordinaire :
    CPU grand public et GPU de 4~Go. Cette contrainte est structurante : elle
    impose un quantization lourde, des prompts courts, une seule génération
    simultanée et un repli CPU~pur configurable.
  \item[N5 --- Latence.] Aucune exigence chiffrée n'avait été fixée au
    démarrage ; la contrainte réelle est venue du matériel. Les mesures de la
    section~\ref{sec:perf} servent de référence, et l'interface est conçue pour
    une latence de plusieurs dizaines de secondes par appel.
  \item[N6 --- Non-régression.] Le Module~1 ne doit subir aucune modification
    fonctionnelle ; le branchement se limite aux quatre fichiers listés en
    section~\ref{sec:branchement}.
  \item[N7 --- Testabilité.] La suite de tests ne doit jamais dépendre
    d'Ollama, d'un téléchargement de poids ni du réseau.
\end{description}

\subsection{Hors périmètre assumé}

Trois points ont été explicitement écartés, et non oubliés :

\begin{itemize}[leftmargin=1.6em,itemsep=1pt,topsep=2pt]
  \item le \emph{streaming} token par token, pour des raisons de format de
    sortie (section~\ref{sec:no-stream}) ;
  \item le \emph{déploiement} : le module suppose un Ollama local, il n'a pas
    été conteneurisé ni publié ;
  \item l'\emph{adaptation à une offre stockée} : le dépôt du Module~1 ne
    conserve aucune offre d'emploi. Le point de branchement est néanmoins en
    place et documenté (\file{cv\_bridge.get\_job\_offer\_context},
    section~\ref{sec:limites}).
\end{itemize}
